\chapter{Capítulo 4: (TODO: título del capítulo)}
\label{cap:capitulo4}

Durante el proyecto se ha desarrollado un sistema de análisis musical bajo una arquitectura secuencial dividida en cinco fases diferenciadas:
\begin{itemize}
\item \textbf{Fase de descubrimiento de archivos.} Inspección recursiva del directorio raíz del corpus para identificar y catalogar archivos con extensiones compatibles (.xml, .musicxml, .mxl, .mei, etc.).
\item \textbf{Fase de preprocesamiento y normalización de archivos.} Limpieza de cadenas de texto para corregir la aparición de, por ejemplo, caracteres no aptos o espacios en los nombres de archivo de las obras, e incompatibilidades de sintaxis entre herramientas como \textit{music21} y los formatos utilizados antes de la llamada al intérprete musical.
\item \textbf{Fase de análisis simbólico y extracción de metadatos.} Carga de la información sobre cada obra en memoria para extraer de ella datos como la tonalidad, el compás, el tempo, la inferencia híbrida de la tonalidad y la detección de patrones rítmicos.
\item \textbf{Fase de análisis semántico de la lírica mediante LLMs locales.} Extracción del texto lírico estructurado de cada obra y su posterior envío a un LLM local para su etiquetado. El modelo determina qué temas trata la letra, asignándole, como máximo, cuatro temáticas a cada canción (Nana, Religión, Naturaleza, Picaresca, Amor, etc.).
\item \textbf{Fase de consolidación, estructuración y almacenamiento en la base de datos SQLite.} Compilación de los diccionarios que contienen las características mencionadas con anterioridad y su exportación a una base de datos SQLite.
\end{itemize}

(TODO: ¿se añade el prompt?)

Para mejorar el rendimiento de la búsqueda de características concretas en el corpus se han diseñado e implementado una serie de funciones que abordan de manera más específica y técnica, basándose en teoría musical, los desafíos de extracción musical y semántica que presentan los motores de búsqueda con modelos de IA generales. Funcionalidades implementadas:
\begin{itemize}
\item \textbf{Robustez en el parseo y extracción híbrida de metadatos.} Dadas las limitaciones de \textit{music21} a la hora de procesar ciertos metadatos estructurados de los archivos, se ha implementado un sistema que reemplaza los atributos de posición de sílaba problemáticos (como \textit{wordpos="s"} por \textit{"i"}) y genera un archivo temporal que se parsea si generar excepciones. Además, si el archivo a procesar es de formato MEI, aparece en el proceso un scraper personalizado, encargado de extraer datos vitales omitidos habitualmente por la importación estándar, como la autora o la armadura de la obra.
\item \textbf{Heurística de inferencia híbrida de tonalidad.} En algunos casos, el analizador de \textit{music21} falla debido a la brevedad o falta de cierta información de las piezas del corpus. Para solucionar este problema se ha implementado un algoritmo híbrido basado en teoría modal e instrumental que analiza toda la obra con el objetivo de determinar su tonalidad. La herramienta pasa por las siguientes fases:
\begin{itemize}
    \item Extrae el número de alteraciones (sostenidos o bemoles) de la armadura de la obra.
    \item Identifica una potencial tónica analizando la última nota de la obra, asumiendo el principio musicológico de la resolución cadencial sobre la tónica.
    \item (TODO: no me gusta cómo hago este cálculo) Si la última nota coincide con una de las dos opciones relativas de la armadura (Mayor o menor), se escoge dicha tonalidad. Si no coincide (por ejemplo, en el caso de obras inconclusas o modales), calcula la frecuencia de aparición de ambas tónicas candidatas a lo largo de la obra, asignando la tonalidad al modo con mayor peso estadístico.
\end{itemize}
\item \textbf{Algoritmo de inferencia de modo musical.} El analizador por defecto de \textit{music21} utiliza el algoritmo de Krumhansl-Schmuckler, que compara la distribución de notas con perfiles de tonalidades mayores y menores. Como no modela explícitamente modos (dórico, frigio, etc.), una pieza modal puede clasificarse erróneamente como la tonalidad mayor o menor más cercana. En este caso, para identificar el modo de la obra, se siguen los siguientes pasos:
\begin{itemize}
    \item Identifica la tónica de la obra. En este caso, se escoge la última nota de la obra, asumiendo que, generalmente, en la música folclórica, la última nota de la melodía es la tónica que resuelve la obra.
    \item Analiza la distancia en semitonos de cada uno de los intervalos con respecto a la tónica y mide cuál de las estructuras modales existentes coincide más con el esquema de intervalos recopilado.
\end{itemize}
\item \textbf{Algoritmos de detección rítmica avanzada.}
    \begin{itemize}
    \item \textbf{Hemiolias horizontales.} Evalúa el mapa global de tiempos de ataque de un compás. Si una pieza en $6/8$ (métrica binaria de subdivisión ternaria) presenta ataques exclusivos en subdivisiones de negra regular (0.0, 1.0, 2.0), se clasifica como una mutación interna a sensación de $3/4$. A la inversa, si un compás de $3/4$ agrupa sus ataques únicamente en 0.0 y 1.5 (dos pulsos de negra con puntillo), se cataloga como estructura hemiólica horizontal binaria.
    \item \textbf{Hemiolias verticales.} Analiza la polifonía o la existencia de múltiples capas internas con distintas distribuciones de ataques en un mismo compás. El algoritmo extrae y compara las matrices de los tiempos de ataque de todas las combinaciones de voces posibles en el compás. Si se detecta que la voz $i$ ejecuta un patrón puramente binario mientras la voz $j$ ejecuta simultáneamente uno ternario, el compás se etiqueta positivamente como hemiola vertical o polirritmia.
    \item \textbf{Síncopas.} El sistema examina la interacción de cada nota con la estructura de acentos del compás. Una síncopa aparece formalmente cuando una nota se inicia en un pulso débil, pero su duración se prolonga lo suficiente para atravesar un acento mayor que en el que comenzó.
    \end{itemize}
\item \textbf{Extracción lírica e integración con LLM local.} Para realizar el análisis poético de temas sobre las obras se ha implementado una función que recorre linealmente la partitura, extrayendo las cadenas de texto asociadas a las sílabas de los versos y consolidándolas en un texto continuo, facilitando así la tarea de análisis del LLM. Posteriormente, el título extraído de la pieza y la letra consolidada se introducen en un $prompt$ estructurado enviado al modelo. El modelo de lenguaje se configura mediante un rol de sistema (experto en musicología y folklore) que lo obliga a devolver exclusivamente un esquema en formato JSON con los temas que describen la obra. Esto permite la automatización del proceso, obteniendo el análisis temático de un LLM dentro de variables limpias y organizadas, sin necesidad de realizar un etiquetado manual.
\item \textbf{Almacenamiento en la base de datos SQLite.} Tras seguir todo el procedimiento de extracción de información, todos los datos se vuelcan en la base de datos pertinente.
\end{itemize}